\section{The appointed texts}
Latin follows the identified 1962 appointments, with the received printed
abbreviations. Biblical study English is Douay--Rheims/Challoner.
Lasance/Walsh supplies historical English for the composed prayers and
the adapted Wisdom chant; Cummiskey supplies the Common Preface.
The shorter historical English conclusion cues are preserved as printed.
They do not expand the longer Latin cues.

\subsection{Introit}
\label{proper-introit}
\properlocus{1962 Missal, printed p.~702 (facsimile p.~783) no.~3974; Mame 1922, p.~748 (facsimile p.~862)}
\begin{latinproper}
Gaudeamus omnes in Domino, diem festum celebrantes sub honore Sanctorum omnium: de quorum solemnitate gaudent Angeli, et collaudant Filium Dei. Ps. 32, 1 Exsultate, iusti, in Domino: rectos decet collaudatio. V. Gloria Patri.
\end{latinproper}
\begin{englishwitness}{Lasance/Walsh, revised 1945}
Let us all rejoice in the Lord, celebrating a feast in honor of all the saints, in whose solemnity the angels rejoice and join in praising the Son of God. Rejoice in the Lord, ye just: praise becometh the upright. V. Glory.
\end{englishwitness}
\properlocus{Historical English source: Lasance/Walsh revised 1945, p.~1265, facsimile p.~1266.}

\subsection{Collect}
\label{proper-collect}
\properlocus{1962 Missal, printed p.~702 (facsimile p.~783) no.~3975; Mame 1922, p.~748 (facsimile p.~862)}
\begin{latinproper}
Omnipotens sempiterne Deus, qui nos omnium Sanctorum tuorum merita sub una tribuisti celebritate venerari: quaesumus; ut desideratam nobis tuae propitiationis abundantiam, multiplicatis intercessoribus, largiaris. Per Dominum.
\end{latinproper}
\begin{englishwitness}{Lasance/Walsh, revised 1945}
O almighty, eternal God, Who hast granted us to venerate the merits of all Thy saints in one solemnity, vouchsafe to us, we beseech Thee, through the multitude of our intercessors, that abundance of Thy mercy for which we ever have hoped. Through our Lord.
\end{englishwitness}
\properlocus{Historical English source: Lasance/Walsh revised 1945, p.~1265, facsimile p.~1266.}

\subsection{Commemoration Collect}
\label{proper-commemoration-collect}
\properlocus{1962 Missal, printed p.~416 (facsimile p.~497) no.~1720; Pustet 1862, p.~357 (facsimile p.~443), petition with Per Dominum cue}
\begin{latinproper}
Absolve, quaesumus, Domine, tuorum delicta populorum: ut a peccatorum nexibus, quae pro nostra fragilitate contraximus, tua benignitate liberemur. Per Dominum nostrum Iesum Christum, Filium tuum: Qui tecum vivit et regnat in unitate.
\end{latinproper}
\begin{englishwitness}{Lasance/Walsh, revised 1945}
Remit, we beseech Thee, O Lord, the sins of Thy people, that by Thy kindness we may be delivered from the trammels of our sins, in which, through our frailty, we have become entangled. Through our Lord.
\end{englishwitness}
\properlocus{Historical English source: Lasance/Walsh revised 1945, p.~733, facsimile p.~731.}

\subsection{Lesson}
\label{proper-lesson}
\properlocus{1962 Missal, printed p.~702–703 (facsimile p.~783–784) no.~3976; Mame 1922, p.~748–749 (facsimile p.~862–863); Apoc.~7:2–12 with liturgical incipit}
\begin{latinproper}
In diebus illis: Ecce ego Ioannes vidi alterum Angelum ascendentem ab ortu solis, habentem signum Dei vivi: et clamavit voce magna quatuor Angelis, quibus datum est nocere terrae et mari, dicens: Nolite nocere terrae et mari, neque arboribus, quoadusque signemus servos Dei nostri in frontibus eorum. Et audivi numerum signatorum, centum quadraginta quatuor milia signati, ex omni tribu filiorum Israel. Ex tribu Iuda duodecim milia signati. Ex tribu Ruben duodecim milia signati. Ex tribu Gad duodecim milia signati. Ex tribu Aser duodecim milia signati. Ex tribu Nephtali duodecim milia signati. Ex tribu Manasse duodecim milia signati. Ex tribu Simeon duodecim milia signati. Ex tribu Levi duodecim milia signati. Ex tribu Issachar duodecim milia signati. Ex tribu Zabulon duodecim milia signati. Ex tribu Ioseph duodecim milia signati. Ex tribu Beniamin duodecim milia signati. Post haec vidi turbam magnam, quam dinumerare nemo poterat, ex omnibus gentibus, et tribubus, et populis, et linguis: stantes ante thronum, et in conspectu Agni, amicti stolis albis, et palmae in manibus eorum: et clamabant voce magna, dicentes: Salus Deo nostro, qui sedet super thronum, et Agno. Et omnes Angeli stabant in circuitu throni, et seniorum, et quatuor animalium: et ceciderunt in conspectu throni in facies suas, et adoraverunt Deum, dicentes: Amen. Benedictio, et claritas, et sapientia, et gratiarum actio, honor, et virtus, et fortitudo Deo nostro in saecula saeculorum. Amen.
\end{latinproper}
\begin{englishwitness}{Douay--Rheims/Challoner study text}
And I saw another angel ascending from the rising of the sun, having the sign of the living God. And he cried with a loud voice to the four angels to whom it was given to hurt the earth and the sea, Saying: Hurt not the earth nor the sea nor the trees, till we sign the servants of our God in their foreheads. And I heard the number of them that were signed. An hundred forty-four thousand were signed, of every tribe of the children of Israel. Of the tribe of Juda, twelve thousand signed: Of the tribe of Ruben, twelve thousand signed: Of the tribe of Gad, twelve thousand signed: Of the tribe of Aser, twelve thousand signed: Of the tribe of Nephthali, twelve thousand signed: Of the tribe of Manasses, twelve thousand signed: Of the tribe of Simeon, twelve thousand signed: Of the tribe of Levi, twelve thousand signed: Of the tribe of Issachar, twelve thousand signed: Of the tribe of Zabulon, twelve thousand signed: Of the tribe of Joseph, twelve thousand signed: Of the tribe of Benjamin, twelve thousand signed. After this, I saw a great multitude, which no man could number, of all nations and tribes and peoples and tongues, standing before the throne and in sight of the Lamb, clothed with white robes, and palms in their hands. And they cried with a loud voice, saying: Salvation to our God, who sitteth upon the throne and to the Lamb. And all the angels stood round about the throne and the ancients and the four living creatures. And they fell down before the throne upon their faces and adored God, Saying: Amen. Benediction and glory and wisdom and thanksgiving, honour and power and strength, to our God, for ever and ever. Amen.
\end{englishwitness}
The Missal adds its liturgical opening. Lasance/Walsh begins:
``In those days, behold, I, John, saw another angel”. The
complete biblical study selection above begins at Apocalypse 7:2.
\properlocus{Historical English source: Lasance/Walsh, p.~1265–1266 (facsimile p.~1266–1267).}

\subsection{Gradual}
\label{proper-gradual}
\properlocus{1962 Missal, printed p.~703 (facsimile p.~784) no.~3977; Mame 1922, p.~749 (facsimile p.~863); Ps.~33:10 and second clause of verse 11}
\begin{latinproper}
Timete Dominum, omnes Sancti eius: quoniam nihil deest timentibus eum. V. Inquirentes autem Dominum, non deficient omni bono.
\end{latinproper}
\begin{englishwitness}{Douay--Rheims/Challoner study text}
Fear the Lord, all ye his saints: for there is no want to them that fear him. but they that seek the Lord shall not be deprived of any good.
\end{englishwitness}
The second sentence is only the latter clause of Psalm 33:11;
its initial lower-case word retains the Bible quotation.

\subsection{Alleluia}
\label{proper-alleluia}
\properlocus{1962 Missal, printed p.~703 (facsimile p.~784) no.~3978; Mame 1922, p.~749 (facsimile p.~863); Matt.~11:28}
\begin{latinproper}
Alleluia, alleluia. V. Venite ad me, omnes qui laboratis et onerati estis: et ego reficiam vos. Alleluia.
\end{latinproper}
\begin{englishwitness}{Douay--Rheims/Challoner study text}
Come to me all you that labor and are burdened, and I will refresh you.
\end{englishwitness}
\begin{englishwitness}{Lasance/Walsh: received chant framing}
Alleluia, alleluia. V. Come to Me, all you that labor and are heavy laden, and I will refresh you. Alleluia.
\end{englishwitness}
\properlocus{Historical English source: Lasance/Walsh, p.~1267 (facsimile p.~1268).}

\subsection{Gospel}
\label{proper-gospel}
\properlocus{1962 Missal, printed p.~703 (facsimile p.~784) no.~3979; Mame 1922, p.~749 (facsimile p.~863); Matt.~5:1–12a; Credo is outside body}
\begin{latinproper}
In illo tempore: Videns Iesus turbas, ascendit in montem, et cum sedisset, accesserunt ad eum discipuli eius, et aperiens os suum, docebat eos, dicens: Beati pauperes spiritu: quoniam ipsorum est regnum caelorum. Beati mites: quoniam ipsi possidebunt terram. Beati qui lugent: quoniam ipsi consolabuntur. Beati qui esuriunt et sitiunt iustitiam: quoniam ipsi saturabuntur. Beati misericordes: quoniam ipsi misericordiam consequentur. Beati mundo corde: quoniam ipsi Deum videbunt. Beati pacifici: quoniam filii Dei vocabuntur. Beati qui persecutionem patiuntur propter iustitiam: quoniam ipsorum est regnum caelorum. Beati estis cum maledixerint vobis, et persecuti vos fuerint, et dixerint omne malum adversum vos, mentientes, propter me: gaudete et exsultate, quoniam merces vestra copiosa est in caelis.
\end{latinproper}
\begin{englishwitness}{Douay--Rheims/Challoner study text}
And seeing the multitudes, he went up into a mountain, and when he was set down, his disciples came unto him. And opening his mouth he taught them, saying: Blessed are the poor in spirit: for theirs is the kingdom of heaven. Blessed are the meek: for they shall possess the land. Blessed are they that mourn: for they shall be comforted. Blessed are they that hunger and thirst after justice: for they shall have their fill. Blessed are the merciful: for they shall obtain mercy. Blessed are the clean of heart: for they shall see God. Blessed are the peacemakers: for they shall be called the children of God. Blessed are they that suffer persecution for justice’ sake: for theirs is the kingdom of heaven. Blessed are ye when they shall revile you, and persecute you, and speak all that is evil against you, untruly, for my sake: Be glad and rejoice for your reward is very great in heaven.
\end{englishwitness}
The received Lasance/Walsh opening is ``At that time, Jesus
seeing the multitudes”. The biblical study selection ends
before the final prophets clause of Matthew 5:12, as the Missal does.
\properlocus{Historical English source: Lasance/Walsh, p.~1267 (facsimile p.~1268).}

\subsection{Offertory}
\label{proper-offertory}
\properlocus{1962 Missal, printed p.~703 (facsimile p.~784) no.~3980; Mame 1922, p.~750 (facsimile p.~864); selected/adapted Wis.~3:1,2,3}
\begin{latinproper}
Iustorum animae in manu Dei sunt, et non tanget illos tormentum malitiae: visi sunt oculis insipientium mori: illi autem sunt in pace, alleluia.
\end{latinproper}
\begin{englishwitness}{Lasance/Walsh, revised 1945}
The souls of the just are in the hand of God, and the torment of malice shall not touch them: in the sight of the unwise they seemed to die, but they are in peace. Alleluia.
\end{englishwitness}
\properlocus{Historical English source: Lasance/Walsh, p.~1267–1268 (facsimile p.~1268–1269).}

\subsection{Secret}
\label{proper-secret}
\properlocus{1962 Missal, printed p.~703 (facsimile p.~784) no.~3981; Mame 1922, p.~750 (facsimile p.~864)}
\begin{latinproper}
Munera tibi, Domine, nostrae devotionis offerimus: quae et pro cunctorum tibi grata sint honore iustorum, et nobis salutaria, te miserante, reddantur. Per Dominum.
\end{latinproper}
\begin{englishwitness}{Lasance/Walsh, revised 1945}
We offer Thee, O Lord, the gifts of our devotion; may they be made pleasing to Thee for the honor of all the just and, by Thy mercy, be profitable to our selves. Through our Lord.
\end{englishwitness}
\properlocus{Historical English source: Lasance/Walsh revised 1945, p.~1268, facsimile p.~1269.}

\subsection{Commemoration Secret}
\label{proper-commemoration-secret}
\properlocus{1962 Missal, printed p.~417 (facsimile p.~498) no.~1726; Pustet 1862, p.~358 (facsimile p.~444), petition with Per Dominum nostrum cue}
\begin{latinproper}
Pro nostrae servitutis augmento sacrificium tibi, Domine, laudis offerimus: ut, quod immeritis contulisti, propitius exsequaris. Per Dominum nostrum Iesum Christum, Filium tuum: Qui tecum vivit et regnat in unitate.
\end{latinproper}
\begin{englishwitness}{Lasance/Walsh, revised 1945}
We offer Thee, O Lord, the sacrifice of praise for the increase of our service, that Thou mayest graciously complete that which, for no merit of ours, Thou hast granted unto us. Through our Lord.
\end{englishwitness}
\properlocus{Historical English source: Lasance/Walsh revised 1945, p.~735-736, facsimile p.~733.}

\subsection{Common Preface}
\label{proper-preface}
\properlocus{1962 Missal, printed p.~252–253 (facsimile p.~333–334) no.~1057; Cummiskey 1861, p.~xxviii (facsimile p.~39), body only; stable dialogue and Sanctus belong to Ordinary}
\begin{latinproper}
Vere dignum et iustum est, aequum et salutare, nos tibi semper et ubique gratias agere: Domine, sancte Pater, omnipotens aeterne Deus: per Christum Dominum nostrum. Per quem maiestatem tuam laudant Angeli, adorant Dominationes, tremunt Potestates. Caeli, caelorumque Virtutes, ac beata Seraphim, socia exsultatione concelebrant. Cum quibus et nostras voces ut admitti iubeas, deprecamur, supplici confessione dicentes:
\end{latinproper}
\begin{englishwitness}{Cummiskey 1861}
IT is truly meet and just, right and available to salvation, that we should always, and in all places, give thanks to thee, O holy Lord, Father Almighty, eternal God. * Thro' Christ our Lord: by whom the angels praise thy majesty, the dominations adore it, the powers tremble before it, the heavens, the heavenly virtues, and blessed seraphim, with common jubilee glorify it. Together with whom we beseech thee, that we may be admitted to join our humble voices, saying:
\end{englishwitness}
\properlocus{Historical English source: Cummiskey 1861, p.~xxviii (facsimile p.~39).}

\subsection{Communion}
\label{proper-communion}
\properlocus{1962 Missal, printed p.~703 (facsimile p.~784) no.~3982; Mame 1922, p.~750 (facsimile p.~864); Matt.~5:8–10}
\begin{latinproper}
Beati mundo corde, quoniam ipsi Deum videbunt: beati pacifici, quoniam filii Dei vocabuntur: beati qui persecutionem patiuntur propter iustitiam, quoniam ipsorum est regnum caelorum.
\end{latinproper}
\begin{englishwitness}{Douay--Rheims/Challoner study text}
Blessed are the clean of heart: for they shall see God. Blessed are the peacemakers: for they shall be called the children of God. Blessed are they that suffer persecution for justice’ sake: for theirs is the kingdom of heaven.
\end{englishwitness}

\subsection{Postcommunion}
\label{proper-postcommunion}
\properlocus{1962 Missal, printed p.~703 (facsimile p.~784) no.~3983; Mame 1922, p.~750 (facsimile p.~864), same petition but Per Dominum cue}
\begin{latinproper}
Da, quaesumus, Domine, fidelibus populis omnium Sanctorum semper veneratione laetari: et eorum perpetua supplicatione muniri. Per Dominum nostrum.
\end{latinproper}
\begin{englishwitness}{Lasance/Walsh, revised 1945}
Grant to Thy faithful people, we beseech Thee, O Lord, ever to rejoice in the veneration of all Thy saints, and to be defended by their unceasing prayers. Through our Lord.
\end{englishwitness}
\properlocus{Historical English source: Lasance/Walsh revised 1945, p.~1268, facsimile p.~1269.}

\subsection{Commemoration Postcommunion}
\label{proper-commemoration-postcommunion}
\properlocus{1962 Missal, printed p.~417 (facsimile p.~498) no.~1728; Pustet 1862, p.~358 (facsimile p.~444), petition with Per Dominum cue}
\begin{latinproper}
Quaesumus, omnipotens Deus: ut, quos divina tribuis participatione gaudere, humanis non sinas subiacere periculis. Per Dominum nostrum Iesum Christum, Filium tuum: Qui tecum vivit et regnat in unitate.
\end{latinproper}
\begin{englishwitness}{Lasance/Walsh, revised 1945}
We pray Thee, O almighty God, that Thou wilt not suffer to be overcome by human dangers those to whom Thou grantest participation in things that are divine. Through our Lord.
\end{englishwitness}
\properlocus{Historical English source: Lasance/Walsh revised 1945, p.~736, facsimile p.~734.}

